\section{Related Work}
\label{sec:related}

\subsection{Rule-Based and Learned Intrusion Detection}
Classic network intrusion detection systems emphasize transparent signatures and protocol analysis. Snort operationalized lightweight signature matching~\cite{roesch1999snort}, while Bro/Zeek demonstrated stateful network analysis and policy-oriented event processing~\cite{paxson1999bro}. These systems remain valuable because their behavior is auditable and their false-positive profile can be tightly controlled for known patterns. Their principal limitation is coverage: a rule cannot detect a semantic variant it does not encode, and authorization or business-logic attacks may depend on context that is not reducible to a payload substring.

Machine learning broadens the decision boundary beyond explicit signatures. DeepLog models execution logs as sequences for anomaly detection~\cite{du2017deeplog}; LUCID demonstrates a lightweight deep-learning approach to DDoS classification~\cite{doriguzzi2020lucid}; and recent BERT-based intrusion detectors apply contextual encoders to network features~\cite{yang2025bertids}. Surveys of explainable intrusion detection note that learned detectors frequently trade interpretability and distributional stability for broader pattern recognition~\cite{neupane2022xids}. Our work differs in both input and output: we analyze request-side HTTP evidence as structured text and require a typed security decision rather than a single anomaly score.

\subsection{Language Models for Structured Security Decisions}
Transformers provide a general sequence modeling foundation~\cite{vaswani2017attention}. Modern instruction models can consume heterogeneous text and produce structured outputs, making them plausible components for security triage, log interpretation, and policy assistance. We use Qwen3-4B-Instruct-2507, a compact dense instruction model in the Qwen3 family~\cite{qwen3,qwen2507}. Rather than evaluate general conversational quality, we study narrow task adaptation under a fixed schema and a request-only evidence boundary.

Parameter-efficient adaptation reduces the cost of specializing a base model. LoRA freezes pretrained weights and learns low-rank updates~\cite{hu2021lora}; QLoRA combines low-rank adapters with a frozen 4-bit quantized base model~\cite{dettmers2023qlora}. These methods make single-GPU specialization feasible, but they do not guarantee reliable uncertainty or domain transfer. Our results show that QLoRA can shift a compact model toward high attack recall while simultaneously eliminating abstention behavior, illustrating why adaptation must be evaluated on out-of-distribution and unknown examples.

\subsection{Prompt Injection and Untrusted Instructions}
Prompt injection attacks exploit the inability of instruction-following models to reliably distinguish trusted control text from attacker-controlled content. PromptInject demonstrated goal hijacking and prompt leaking through adversarial natural-language inputs~\cite{perez2022promptinject}. Indirect prompt injection extends this threat to content retrieved from external sources or tools~\cite{greshake2023indirect}. HackAPrompt provided large-scale empirical evidence that diverse human-authored prompts can bypass model instructions~\cite{schulhoff2023hackaprompt}. For an HTTP security classifier, the request body itself is untrusted data and may contain text that tells the classifier to ignore its system policy. We therefore treat prompt-injection obedience, raw-output preservation, and message serialization as first-class evaluation concerns.

\subsection{Evaluation, Abstention, and Reproducibility}
Security model evaluation is particularly susceptible to leakage and selective reporting. We use stratified bootstrap confidence intervals~\cite{efron1993bootstrap}, paired McNemar tests~\cite{mcnemar1947note}, and Holm--Bonferroni correction~\cite{holm1979sequentially} to distinguish observed differences from sampling variability. In addition to accuracy and F1, we report safe false-positive rate and unknown-class behavior. The unknown class is operationally analogous to selective prediction or abstention: it represents cases in which model-visible request evidence is insufficient for a justified safe/attack decision. Unlike work that reports only accepted predictions, we score every sample and count schema failures and forced classifications explicitly.

The closest conceptual distinction in our study is between \emph{security judgment} and \emph{output alignment}. A model can fail because it chooses the wrong verdict, because it follows attacker instructions, or because it fails to serialize a valid schema. We report these failure channels separately and compare multiple inference artifacts to avoid attributing every difference to model weights alone.
